\section{MCPと高電圧パルス駆動回路}

\subsection{MCP仕様}
\begin{itemize}
    \item 型番: Hamamatsu F2222-21P-Y010 (2段式)
    \item 特性: 高ゲインが得られる2段式（Chevron）を採用。
\end{itemize}

\subsection{パルス駆動の必要性}
MCPには DC $+4\,\mathrm{kV}$ 程度の高電圧が必要だが、実験開始時（コイル点弧時）の激しい磁場変動時に電圧がかかっていると、誘導ノイズで画像が乱れるだけでなく、MCPが焼損するリスクがある。
そのため、計測タイミング（数ms間）のみ高電圧を印加する「ゲーティング回路」を自作して使用する\cite{Takeda2023}。

\subsection{パルス回路構成と主要部品}
本装置で使用している回路は、以下の部品で構成される。

\begin{description}
    \item[高耐圧MOSFET] スイッチング素子。
    \begin{itemize}
        \item 型番: \textbf{IXYS IXTT02N450HV}
        \item 定格: 耐圧 $4500\,\mathrm{V}$ (これ以下のFETは耐圧不足で即死するため注意)。
    \end{itemize}
    
    \item[ゲートドライバ (フォトカプラ)] 高電圧側と制御側を絶縁する。
    \begin{itemize}
        \item 型番: \textbf{RENESAS PS9531}
        \item 特徴: 高速スイッチングが可能。
    \end{itemize}
    
    \item[絶縁DC-DCコンバータ] フローティング電源（ゲート駆動用15V）を作る。
    \begin{itemize}
        \item 型番: \textbf{TRACO POWER TIM 2-0923}
    \end{itemize}
    
    \item[高圧電源 (DC)] ベースとなる電源。
    \begin{itemize}
        \item Phosphor用: 松定プレシジョン HPMR-5P (+5kVまで)
        \item MCP-in用: 松定プレシジョン HPMR-2N (-2kVまで)
    \end{itemize}
\end{description}

\subsection{結線上の注意}
MOSFETはPhosphor側（正電位）とMCP-in側（負電位）の両方に入れ、同時タイミングでゲートを開くことで、MCPにかかる電位差を一瞬で形成する設計となっている。